\documentclass[11pt]{article}
\usepackage[T1]{fontenc}
\usepackage{lmodern,amsmath,amssymb,amsthm,geometry,hyperref}
\geometry{margin=25mm}
\newtheorem{theorem}{Theorem}
\newtheorem{lemma}{Lemma}
\newcommand{\PP}{\mathcal P}
\newcommand{\MM}{\mathfrak M}
\newcommand{\BB}{\mathcal B}
\title{Outer arithmetic value on the original quotient:\\
independent reconstruction and finite estimates}
\date{19 September 2026}
\begin{document}
\maketitle

This receiving proof concerns the outer part of the arithmetic density.
The original source, packet, relation ideal, four endpoint signs and
Gamma mass are retained.  The original EIQ equilibrium and RTM
one-factor bounds are the proved programme providers specified in the
accompanying acceptance record; no arithmetic allocation or projected
sign is inferred from a full-quotient determinant.

\section{Objects and exact receiver}
Fix the original stipulated quartet
\(\rho=1/2+\delta+i\gamma\), \(0<\delta<1/2\), \(\gamma>2\),
with original multiplicity \(m\geq1\).  Define
\[
 h(s)=\prod_{z\in\{\rho,\bar\rho,1-\rho,1-\bar\rho\}}(s-z)^m,\qquad
 w_h(y)=\frac{|(2\xi/h)(1/2+iy)|^2}{2\pi},\qquad m_k=w_h^{*k}.
\]
Keep \(k=4\ell+1\), \(e=1+k(m-1)\), \(q=e(k+1)^2\), \(c=k/2\), and
\[
 \chi_k(S)=
 \prod_{a,b_0=0}^k
 [S-c-(2a-k)\delta-i(2b_0-k)\gamma]^e,\qquad
 Q(y)=i^{-q}\chi_k(c+iy).
\]
Thus \(Q\) is real, even and monic of degree \(q\); all its roots,
including every repetition, have modulus at most
\(R=k(\delta^2+\gamma^2)^{1/2}\).
Write
\[
 \sigma(y)=\frac{|\Gamma(1/4+iy/2)|^2}{2\pi},\quad
 \int_{\mathbb R}\sigma=\sqrt{2\pi},\quad
 \alpha=\pi/2,\quad
 M(\theta)=\int e^{\theta y}w_h(y)\,dy,\quad
 \MM=M(\alpha),\quad b=(\log M)'(\alpha)>0 .
\]
For a positive density \(\nu\), \(H_N[\nu]\) is the Gram of
\(1,y,\ldots,y^N\), \(J_N\) is remainder modulo the full \(Q\), and
\[
 G_N[\nu]=(J_NH_N[\nu]^{-1}J_N^*)^{-1},\qquad
 \BB[\nu]=\log\det G_{q-1}[\nu]+\log\det G_q[\nu]
 -\log\det G_{2q-1}[\nu]-\log\det G_{2q}[\nu].
 \tag{OR1}
\]
The monic frame consisting of the first \(q\) monomials followed by
\(Q,yQ,\ldots,y^{N-q}Q\) has determinant one.  Its Schur complement
proves, with its cross blocks retained,
\[
 \det G_N[\nu]=\frac{\det H_N[\nu]}
 {\det\operatorname{Gram}_{\nu}(Q,yQ,\ldots,y^{N-q}Q)}. \tag{OR2}
\]
The polynomial substitution \(P(S)\mapsto P(c+iy)\) is invertible, has
coefficient determinant of modulus one, and carries \(\chi_kP\) to
\(i^q QP(c+iy)\); this proves that (OR1)--(OR2) compute the original
quotient.  A positive scalar \(A\) multiplies every \(G_N\) by \(A\).
Each quotient has dimension \(q\), hence \(\BB[A\nu]=\BB[\nu]\).

Set
\[
 d=4m-2,\quad p=2d-\tfrac12,\quad a=p+1,\qquad
 \tau_d=(1+y^2)^{-d}\sigma,\qquad
 \omega_k(y)=
 \begin{cases}\sigma(y),&|y|\leq kb,\\
 \MM^{-k}m_k(y),&|y|>kb.\end{cases}                    \tag{OR3}
\]
The notation \(a\) in (OR3) is a tail exponent, distinct from the
integer indexing a root in \(\chi_k\).

\begin{theorem}
On the original fixed-packet cofinal range, with all original finite
guards retained,
\[
 \BB[\omega_k]-\BB[\sigma]
 =-d\mathcal C_\Gamma q+\mathcal E_k,\qquad
 \mathcal C_\Gamma=2\mathcal L(2,\pi)-8\log2+4,          \tag{OR4}
\]
where \(\mathcal L(\alpha_E,\beta)=\int\log x\,\rho_{\alpha_E,\beta}(x)\,dx\)
is the original EIQ logarithmic moment and
\[
 |\mathcal E_k|\leq C_h\left[
 \frac{q\log\log(q+16)}{\log(q+2)}
 +q^{199/200}\log^2(q+2)\right]=o_h(q).                \tag{OR5}
\]
In particular the outer term has no nonzero contribution of magnitude
\(q\log k\) or \(q\log\log k\).  This theorem is about the outer density
in (OR3); it does not delete the central density in \(m_k\).
\end{theorem}

\section{The growing convolution, including both tails}
The exact identity following from the definition of the convolution is
\[
 f(y)=\MM^{-1}e^{\alpha y}w_h(y),\qquad
 \MM^{-k}m_k(y)=e^{-\alpha y}f^{*k}(y).                \tag{OR6}
\]
The auxiliary \(f\) has mass one, mean \(b\) and variance \(V>0\).
The original density is recovered with the entire factor in (OR6).
From the completed-zeta definition one obtains directly
\[
 w_h(y)=\frac{(y^2+1/4)^2}{\sqrt\pi}\,
 \frac{|\zeta(1/2+iy)|^2}
 {[(y-\gamma)^2+\delta^2]^{2m}[(y+\gamma)^2+\delta^2]^{2m}}\sigma(y).
\]
Stirling's formula \cite{DLMFgamma}, on this fixed vertical line, gives
\[
 f(y)=c_fy^{-a}Z(y)(1+O_h(y^{-1})),\quad
 c_f=\frac{\sqrt{2/\pi}}{\MM},\quad
 Z(y)=|\zeta(1/2+iy)|^2,\quad y\longrightarrow+\infty. \tag{OR7}
\]
The RTM envelope gives \(f(y)\leq C_h(1+|y|)^{-p}\), \(p\geq7/2\);
on the negative half-line it has an additional
\(e^{-2\alpha|y|}\) factor.  Thus its \(9/4\)-moment is finite.
The mean-square zeta formula gives, whenever \(|z|\leq T^{19/20}\),
\[
 \int_T^{2T}f(y-z)\,dy\leq C_h T^{1-a}\log T.          \tag{OR8}
\]
Indeed \(y-z\) lies in a fixed enlargement of \([T,2T]\), the power
factor in (OR7) is \(O(T^{-a})\), and the zeta integral over that
enlargement is \(O(T\log T)\).

The proved Fourier local-limit estimate in the supplied localization
provider is uniform in the evaluation point:
\[
 \sup_x\left|f^{*n}(nb+x)
 -(2\pi nV)^{-1/2}e^{-x^2/(2nV)}\right|
 \leq C_h n^{-5/8}.                                  \tag{OR9}
\]
Its Fourier proof bounds an \(L^1\) Fourier difference; the factor
introduced by translation has modulus one.  Integrating (OR9) on
\(|x|\leq n^{1/2+1/24}\) costs \(O(n^{-1/12})\).
Chebyshev's inequality for both probability densities bounds the two
complementary masses by \(O(n^{-1/12})\).  If \(g_n\) is the Gaussian
centered at \(nb\), this proves
\[
 \|f^{*n}-g_n\|_1\leq C_h n^{-1/12}.                  \tag{OR10}
\]

For \(q^{9/10}\leq T\leq128q\), the exact definition of \(q\) implies
\[
 c_hT^{1/3}\leq k\leq C_hT^{5/9}.                    \tag{OR11}
\]
Put \(A=T^{19/20}\), \(f_s=f1_{\{|y|\leq A\}}\) and \(f_l=f-f_s\).
Neither subprobability density is normalized.  In the binomial
convolution expansion, terms with at least two large summands have
total mass at most \(C_hk^2A^{2-2p}\).  Dividing by
\(L_T=kT^{-p}\log T\) gives
\[
 C_h kT^pA^{2-2p}/\log T
 \leq C_hT^{-25/36}/\log T.                           \tag{OR12}
\]
The exponent is largest at \(p=7/2\); larger original multiplicity
only improves it.

For \(\lambda=(\log T)/A\), Taylor's formula on the truncated support
gives
\[
 \int e^{\pm\lambda y}f_s(y)\,dy
 \leq\exp(C_h\lambda+C_h\lambda^2T).
\]
Consequently the mass of \(f_s^{*j}\), \(j\leq k\), outside
\([-T/2,T/2]\) is at most
\(2\exp[-c_hT^{1/20}\log T]\), eventually.  Here
\(k/T\to0\) and \(k\lambda\to0\) by (OR11), so the two positive
exponent terms are strictly smaller than the Chernoff exponent.
This bounds the all-small convolution.

The remaining convolution is \(k f_l*f_s^{*(k-1)}\).
On \(A<|z|\leq T/2\), \(f_l(y-z)\leq C_hT^{-p}\), whereas the
rest density has mass at most \(C_h k/A^2\) by its variance and
\(A\geq2(k-1)b\).  Its window integral is at most
\(C_hk^2T^{1-p}/A^2\); divided by \(L_T\), it is
\(O_h(kT/A^2)\), a fixed negative power of \(T\).
On \(|z|>T/2\) use the preceding Chernoff bound and the full mass
of \(f_l\).  On \(|z|\leq A\), one has \(y-z>A\).
Replacing the small rest convolution by the full one has total
variation at most \(C_hkA^{1-p}\).  The window estimate (OR8) bounds
its contribution by \(L_TO_h(kA^{1-p})\).
Replacing that full rest convolution by \(g_{k-1}\), on the same
region, costs \(L_TO_h(k^{-1/12})\) by (OR8) and (OR10).
The complementary Gaussian region is exponentially small, including
the part near \(z=y\), since its center is \(O_h(k)\) and its
variance is \(O_h(k)\).
Finally (OR7) and \(|z|\leq A\) permit replacing
\((y-z)^{-a}\) by \(y^{-a}\), with relative integrated error
\(O_h(A/T)\).  Thus, with \(n=k-1\) and \(G=\sqrt{2nV}\),
\[
 \int_T^{2T}\left|f^{*k}(y)-kc_fy^{-a}\mathcal Z_G(y-nb)\right|dy
 \leq C_h L_T(k^{-1/12}+T^{-1/20}),\qquad
 \mathcal Z_G(t)=\frac1{\sqrt\pi G}\int e^{-(u/G)^2}Z(t+u)\,du .
 \tag{OR13}
\]
This estimate uses an averaged zeta norm; it asserts no pointwise
regular variation across zeta zeros.

\section{The arithmetic smoothing error}
Let \(c_\zeta=2\gamma_{\rm E}-\log(2\pi)\) and
\(\mathcal L_\zeta(t)=\log t+c_\zeta\).  The established theorem
of Ivi\'c \cite[Section 2, (2.1) and its stated \(E\)-analogue]{Ivic}
gives, for \(1\ll U\leq\sqrt T/2\),
\[
 \int_T^{2T}|E(t+U)-E(t)|^2dt
 \ll_\varepsilon TU\log^3 T+
 T^{1/2+\varepsilon}U^2+T^{1+\varepsilon}U^{1/2},       \tag{OR14}
\]
where \(E(t)=\int_0^tZ(v)\,dv-t(\log(t/(2\pi))+2\gamma_{\rm E}-1)\).
The cited result is a second-moment theorem, without a pointwise or
fourth-moment conjecture.
After truncation at \(|u|\leq G\log T\), integration by parts and
subtraction of \(E(t)\) under the odd kernel give
\[
 \mathcal Z_G(t)-\mathcal L_\zeta(t)
 =\frac2{\sqrt\pi G^3}\int
 u e^{-(u/G)^2}(E(t+u)-E(t))\,du+O(G/T)+O(T^{-100}).
 \tag{OR15}
\]
Polynomial bounds from (OR7) justify the Gaussian tails.  Minkowski's
inequality applied to the square roots of the three terms in
(OR14) yields respectively
\[
 G^{-1/2}\log^{3/2}T,\quad
 T^{-1/4+\varepsilon},\quad T^\varepsilon G^{-3/4}
\]
after division by \(\sqrt T\).  The Gaussian integrals giving these
powers are \(G^{-3}\int_0^\infty u^{3/2}e^{-(u/G)^2}du\),
\(G^{-3}\int_0^\infty u^2e^{-(u/G)^2}du\), and
\(G^{-3}\int_0^\infty u^{5/4}e^{-(u/G)^2}du\).
Negative \(u\) is reduced by translation and a bounded dyadic cover.
For \(|u|\leq1\), positivity of \(Z\) bounds the error increment by
the length-four increment plus \(O(\log T)\); (OR14) at length four
gives the smaller term \(O(T^\varepsilon G^{-3})\).
Since
\(c_hT^{1/6}\leq G\leq C_hT^{5/18}\), all uses of (OR14) lie in
its stated range eventually.  Taking \(\varepsilon=1/1000\),
Cauchy--Schwarz, (OR11)--(OR15), and the exact (OR6) imply
\[
 \int_T^{2T}\left|
 \frac{\MM^{-k}m_k(y)}
 { [k/(\sqrt\pi\MM)]\mathcal L_\zeta(y)(1+y^2)^{-d}\sigma(y)}
 -1\right|dy\leq C_hT^{199/200}.                      \tag{OR16}
\]
The shift \(nb=O_h(k)\) changes the logarithmic main term by
\(O_h(k/T)\).  Division in (OR16) is legitimate because its
power-log denominator on this window is comparable to
\(kT^{-a}\log T\); the exponential factors cancel exactly.
The original evenness of \(m_k,\sigma,\tau_d\) supplies the negative
half-line.

\section{Full polynomial metrics and their attained quotients}
For the Gamma-orthonormal polynomials \(\varphi_j\), define
\(K_N^\sigma(y,y)=\sum_{j=0}^N|\varphi_j(y)|^2\).
The normalization of Koelink--Van der Jeugt
\cite[Proposition 2.1, (2.9)]{KVJ} is transported by
\[
 x=y/2,\quad\lambda=1/4,\quad\phi=\pi/2,\qquad
 \varphi_j(y)=2^{\lambda-1/2}p_j^{(\lambda)}(y/2;\pi/2).
\]
The weighted diagonal density in \(y\) is one half its \(x\) version.
Euler's integral \cite[(15.6.1)]{DLMFhyper}, with its
\(\Gamma(2\lambda)\) conversion from the regularized
hypergeometric function, cancels the weight's
\(|\Gamma(\lambda+ix)|^2\).  Its absolute integrand is
\[
 u^{\lambda-1}(1-u)^{\lambda-1}(1+R_tu)^{-\lambda},
 \qquad R_t=4t/(1-t)^2 .
 \]
For \(t\geq1/2\), splitting at \(1/R_t\) and \(1/2\) bounds this
integral by \(R_t^{-1/4}(16+2\log R_t)\).
Taking \(t=(N+1)/(N+2)\), using positivity of each diagonal
summand and \(t^{-N}\leq e\), proves
\[
 K_N^\sigma(y,y)\sigma(y)\leq20\log(N+2).              \tag{OR17}
\]
For any subspace \(S\subset\PP_{2q}\), orthogonal projection gives
\(K_S\leq K_{2q}^\sigma\) and \(\int K_S\,d\sigma=\dim S\).

In a Gamma-orthonormal frame of \(S\), expansion of the squared
determinant proves
\[
 \det H_S[r\sigma]
 =\frac1{s!}\int|\det[\varphi_i(y_j)]|^2
              \prod_{j=1}^s r(y_j)\,d\sigma(y_j).
 \]
The unweighted right-hand measure has mass one and one-point
density \(K_S\,d\sigma\).  Jensen's inequality and
\(\log\lambda\leq\lambda-1\), applied to the Gram eigenvalues, give
\[
 \int\log r\,K_S\,d\sigma
 \leq\log\frac{\det H_S[r\sigma]}{\det H_S[\sigma]}
 \leq\int(r-1)K_S\,d\sigma.                          \tag{OR18}
\]
This is a statement about the full Gram determinant, including all
cross blocks.

Let \(K=q^{9/10}\) and
\[
 A_k=\frac{k\mathcal L_\zeta(q)}{\sqrt\pi\MM},\qquad
 \widehat\omega_k=
 \begin{cases}A_k\tau_d,&|y|\leq K,\\
 \omega_k,&|y|>K,\end{cases}\qquad
 r_k=\widehat\omega_k/(A_k\tau_d).                    \tag{OR19}
\]
The proved localization bounds, obtained by placing the
\((k-3)\)-fold tilted convolution in its central Gaussian interval
and retaining the three-factor RTM lower tail, give
\[
 c_h k^{-1/2}(1+y^2)^{-M_*}
 \leq\omega_k(y)/\sigma(y)\leq C_hk^{p+1}.            \tag{OR20}
\]
They start at \(|y|=kb\); inside, the quotient is exactly one.
Here \(M_*=21+4m\).  In detail, the interval
\([(k-3)b-\sqrt{(k-3)V},(k-3)b]\) has tilted probability bounded
below by a positive fixed constant, by (OR9).  For \(y\geq kb\)
the remaining positive argument is at most \(2y\), eventually.
Inserting the original three-factor estimate
\(w_h^{*3}(u)\geq c_he^{-\alpha|u|}(1+|u|)^{-42-8m}\)
in that part of the exact convolution proves (OR20); evenness
handles the negative side.  This derivation retains \(\MM^{k-3}\),
which becomes the fixed factor \(\MM^{-3}\) after (OR6).

On \(\PP_{2q}\), Gamma multiplication has norm at most \(2(2q+1)\).
Jensen's inequality applied to the convex function
\((1+x)^{-M_*}\), in the probability
\(|P|^2d\sigma/\|P\|_\sigma^2\), and (OR20), therefore prove
\[
 q^{-C_h}\|P\|_\sigma^2\leq
 \|P\|_{\omega_k}^2,\ \|P\|_{\widehat\omega_k}^2,\
 A_k\|P\|_{\tau_d}^2
 \leq q^{C_h}\|P\|_\sigma^2 .                        \tag{OR21}
\]
The constants here can be enlarged to absorb fixed finite source
constants; the assertion applies on the explicitly cofinal original
sequence.  The same inequality holds for every restriction, and
for a quotient by taking the minimum in each unchanged affine
fibre.

For completeness the small central replacement in (OR19) has its
own finite metric comparison.  The exact generating function
\[
 \sum_{j=0}^\infty\frac{b_j(z)}{j!}t^j
 =(1+t^2)^{-1/4}\exp(z\arctan t),\qquad
 \gamma_j=\|b_j\|_\sigma^2=\sqrt{2\pi}(2j)!/4^j
\]
follows by summing the original three-term recurrence.
Cauchy's coefficient estimate at \(|t|=(2q+1)/(2q+2)\)
and then at \(|z|=2K\) shows, for any \(P\) vanishing to order \(r\),
\[
 \int_{-K}^K|P|^2d\sigma
 \leq A_{q,K,r}\|P\|_\sigma^2,\quad
 A_{q,K,r}=4^{1-r}e^2(2q+1)^2\sqrt{2q+2}
                    [2(2q+2)]^{2K}.                \tag{OR22}
\]
Choose the actual integer \(r\) so that
\(A_{q,K,r}q^{2C_h}\leq q^{-20}\); one may take
\(r=\lceil C'_h(K+1)\log(q+2)\rceil\).
On any actual source or relation space use its own restricted
Taylor map to \(\mathbb C^r\), its full kernel and its attained
image.  The kernel metric changes by relative \(O(q^{-20})\);
its log determinant changes by \(O(q^{-19})\).
The attained image has dimension at most \(r\), so (OR21) costs
at most \(C_h r\log q\).  Factoring the original Gram by this
exact short sequence proves
\[
 \left|\log\frac{\det H_S[\widehat\omega_k]}
                        {\det H_S[\omega_k]}\right|
 \leq C_h q^{9/10}\log^2(q+2)                       \tag{OR23}
\]
for each of the original source and relation spaces.
No quotient frame is discarded; its fixed determinant cancels in
the ratio.

On each dyadic interval \(K\leq |y|\leq128q\), (OR16) bounds
\(r_k-a_q\) in \(L^1\), where
\(a_q(y)=\mathcal L_\zeta(|y|)/\mathcal L_\zeta(q)\).
This \(a_q\) is bounded above and below by fixed positive constants.
The pointwise polynomial bounds in (OR20) imply
\[
 |\log r_k-\log a_q|\leq C_h\log(q+2)|r_k-a_q|.
 \tag{OR24}
\]
To verify this inequality, when \(|r_k-a_q|\leq a_q/2\) use the
bounded derivative of \(\log\) on that interval.  On the complement
the right-hand difference is bounded below by a fixed positive
constant and the logarithms are \(O_h(\log q)\).
Thus (OR17) bounds both error integrals in (OR18) by
\(C_hq^{199/200}\log^2(q+2)\).
For \(q/\log^4q\leq|y|\leq128q\), the deterministic \(a_q-1\)
and its logarithm are \(O(\log\log q/\log q)\); use the exact
kernel mass at most \(2q+1\).  On
\(K\leq|y|\leq q/\log^4q\), use (OR17) and the interval length.
The full remaining tails are bounded by repeated Gamma
multiplication: for \(\deg P\leq2q\),
\[
 \int_{|y|>128q} y^{2d}|P|^2d\sigma
 \leq (128q)^{-2q}[2(3q+d+1)]^{2(q+d)}\|P\|_\sigma^2 .
 \tag{OR25}
\]
Together with (OR20), this gives an exponentially small tail for
both \(r_k-1\) and \(\log r_k\), after summing a Gamma-orthonormal
frame.  Both integration regions are retained.
Writing
\[
 E_q=C_h\left[q^{199/200}\log^2(q+2)
       +q\log\log(q+16)/\log(q+2)\right],
\]
we have therefore proved, uniformly in the required subspaces,
\[
 \int\bigl(|r_k-1|+|\log r_k|\bigr)K_S\,d\sigma\leq E_q.
 \tag{OR26}
\]

The desired reference is \(\tau_d\), and its exact transport is
\(P\mapsto(y+i)^dP\).  It is an isometry from
\(\PP_{N-d}\) in \(\sigma\) onto
\((y+i)^d\PP_{N-d}\) in \(\tau_d\), since
\(|y+i|^{2d}=(1+y^2)^d\).
For the factor \(r_k\) the same identity holds.  On a relation
space the map is \(QP\mapsto Q(y+i)^dP\), preserving the complete
factor \(Q\).  Its omitted quotient is the actual image of
the first \(d\) divided derivatives of \(P\) at \(-i\), hence has
dimension at most \(d\).  Its two metrics have logarithmic
relative cost \(O_h(d\log q)\) by (OR21).
Apply (OR18),(OR26) on the exact image, then include this quotient.
Sum the original signs using (OR2), undo (OR23), and cancel \(A_k\)
only after each quotient's dimension has been retained.  This proves
\[
 \BB[\omega_k]-\BB[\sigma]
 =\BB[\tau_d]-\BB[\sigma]+O_h(E_q).                  \tag{OR27}
\]

\section{The rational model and its finite free-energy error}
Let \(\rho_{\alpha_E,\beta}\) be the original EIQ density for
\(V(x)=\beta\sqrt x-\alpha_E\log x\).  Its Euler inequality is
\(2U-V\leq-\lambda\), and
\(\mathcal F=-\lambda-\iint\log|x-y|\rho(x)\rho(y)\,dxdy\).
The EIQ endpoint formulas give smooth dependence on both positive
parameters and support in a fixed positive compact interval on
the parameter set used below.  The density is bounded there,
has square-root edges, finite logarithmic energy and finite
entropy relative to
\(d\omega(x)=\tfrac12x^{-1/2}e^{-\sqrt x}dx\).

Here is an independent finite-configuration estimate underlying the
shrinking secant.  Replace each configuration point \(x_j\) by
uniform mass on the complex circle \(x_j+\epsilon e^{i\theta}\);
denote the average of these measures by \(\nu_\epsilon\).
The circle mean identity is
\[
 \frac1{2\pi}\int_0^{2\pi}\log|z-\epsilon e^{i\theta}|\,d\theta
 =\log\max(|z|,\epsilon).
 \tag{OR28}
\]
It follows from harmonic mean value outside and inside the circle,
with continuity at its boundary.  For compactly supported
zero-mass signed measures with finite logarithmic energy,
\[
 -\iint\log|z-w|\,d\eta(z)d\eta(w)
 =\frac12\int_0^\infty
 \iint e^{-t|z-w|^2}\,d\eta(z)d\eta(w)\frac{dt}{t}\geq0.
 \tag{OR29}
\]
The scalar logarithm identity proves the equality after truncation;
the zero mass cancels the constant.  The Fourier transform of
the planar Gaussian proves the nonnegative sign.  Absolute
logarithmic integrability of the present circle measures and
bounded density justifies passage to the limit.
Use (OR29) with \(\eta=\nu_\epsilon-\rho\,dx\).
The diagonal circle energy is \(\log\epsilon\), off-diagonal
circle means are at least \(\log|x_i-x_j|\), and (OR28) bounds
the increase of the cross energy by \(2\|\rho\|_\infty\epsilon\).
The Euler inequality now yields
\[
 2\sum_{i<j}\log|x_i-x_j|-n\sum_j V(x_j)
 \leq n^2\mathcal F+n\log n+4n\|\rho\|_\infty
 \quad(\epsilon=1/n).                               \tag{OR30}
\]

The literal parity coordinate is \(y^2=q^2x\).  With
\(B_q(x)=2qe^{\sqrt x}\sigma(q\sqrt x)\), its exact ensemble is
\[
 J_{n,q}(n\alpha_E)=\frac1{n!}\int
 \prod_{i<j}(x_i-x_j)^2
 \prod_j x_j^{n\alpha_E}B_q(x_j)\,d\omega(x_j).
 \tag{OR31}
\]
The original Gamma bound gives
\(B_q(x)\leq Cq^{3/2}
e^{-(\pi q/2-3/2)\sqrt x}\).
Apply (OR30) with
\(\beta_n=\pi q/(2n)-3/(2n)\), and integrate against the
probability \(\omega^n\).
The explicit prefactor and \(n!\) cost \(O(n\log(q+n))\).
Changing \(\beta_n\) to \(\pi q/(2n)\) costs \(O(n)\), since
\(\partial_\beta\mathcal F=-\int\sqrt x\,\rho\,dx\) is uniformly
bounded here.  For the lower bound restrict to the actual
equilibrium support, use the original Gamma lower bound there,
and change measure to \(\rho^n\).  Jensen's inequality gives
the same \(n^2\mathcal F\), with \(O(n)\) missing-diagonal
energy and entropy, plus \(O(n\log(q+n))\) prefactors.  Hence
\[
 |\log J_{n,q}(n\alpha_E)
   -n^2\mathcal F(\alpha_E,\pi q/(2n))|
 \leq Cn\log(q+n),                                  \tag{OR32}
\]
uniformly when \(q/n\) is near \(2\), and \(\alpha_E\) is near
\([7/4,9/4]\).

For the original relation rank \(r=q,q+1\), define
\[
 D_{k,r}(t)=\det\operatorname{Gram}_{(1+y^2)^{qt}\sigma}
                  (Q,yQ,\ldots,y^{r-1}Q),\qquad
 F_{k,r}(t)=q^{-2}\{\log[D_{k,r}(t)/D_{k,r}(0)]
                      -2qrt\log q\}.                \tag{OR33}
\]
The squared-Vandermonde integral proves convexity by
differentiating twice: the second derivative of its logarithm
is a nonnegative variance.  The paired-root comparison on
\(|y|\geq2^{-16}q\), and the complete inner-interval estimate
in the accompanying root-comparison proof, give logarithmic
quadratic-form width \(O_h(1)\) uniformly for \(|t|\leq1/8\).
Its determinant cost is \(O_h(q)\), because the rank is \(r\).
For the reference power \(|y|^{2q(1+t)}\sigma\), its even and
odd blocks have sizes
\[
 n_0=\lceil r/2\rceil,\quad n_1=\lfloor r/2\rfloor,
 \qquad a_0=q(1+t),\quad a_1=q(1+t)+1.
\]
Their parameters satisfy
\[
 (a_i/n_i,\pi q/(2n_i))=(2+2t,\pi)+O(1/q)
 \quad(i=0,1).
 \tag{OR34}
\]
The exact scaling factor of each block is
\(q^{2a_i n_i+2n_i(n_i-1)}\).
Its change is exactly \(2qrt\log q\) after the two blocks
are added.  Applying (OR32) and EIQ smoothness therefore gives
\[
 \sup_{|t|\leq1/8}|F_{k,r}(t)-F_*(t)|
 \leq C_h\frac{\log(q+2)}q,\qquad
 F_*(t)=\tfrac12[\mathcal F(2+2t,\pi)-\mathcal F(2,\pi)] .
 \tag{OR35}
\]
Here \(F_*'(0)=\mathcal L(2,\pi)\) and \(F_*''\) is bounded.
Take \(\epsilon=\sqrt{\log(q+2)/q}\) with
\(d/q<\epsilon/2\) and \(2\epsilon<1/8\).
The convex secant over \([-d/q,0]\) is bounded below and
above by adjacent secants of length \(\epsilon\).
Equation (OR35) changes these latter slopes by at most
\(C_h\log(q+2)/(q\epsilon)\); bounded \(F_*''\) changes their
centers from \(F_*'(0)\) by \(O_h(\epsilon)\).
Returning the exact powers in (OR33) proves
\[
 \log\frac{D_{k,r}(-d/q)}{D_{k,r}(0)}
 =-2dr\log q-dq\mathcal L(2,\pi)
   +O_h(\sqrt{q\log(q+2)}),\qquad r=q,q+1.            \tag{OR36}
\]
No derivative of an asymptotic remainder was taken.

For the source determinant, use the monic frame
\[
 1,\ldots,y^{d-1},(y+i)^d,\ldots,y^{N-d}(y+i)^d .
\]
Its determinant is one and the last block is exactly the Gamma
Gram on \(\PP_{N-d}\).
The first \(d\) divided derivatives at \(-i\) give the quotient
and have determinant one on the first \(d\) monomials.
Their minimum Gram is bounded above by the fixed monomial
trial lifts.  It is bounded below by a negative fixed power
of \(N+2\): use
\(\tau_d\succeq[1+4(N+1)^2]^{-d}\sigma\) on \(\PP_N\),
then the above generating function on fixed circles about
\(-i\) bounds every Gamma jet covariance by a fixed power
of \(N+2\).  There are only \(d\) rows, so the reciprocal
covariance lower bound is again polynomial.
Consequently, retaining this quotient determinant,
\[
 \log\frac{\det H_N[\tau_d]}{\det H_N[\sigma]}
 =-\sum_{j=N-d+1}^N\log\gamma_j+O_d(\log(N+2)).       \tag{OR37}
\]
The exact \(\gamma_j=\sqrt{2\pi}(2j)!/4^j\) gives
\(\log\gamma_j=2j\log j-2j+O(\log(j+2))\).
Insert the four source endpoints and their original signs:
\[
 \sum_{N=q-1,q,2q-1,2q}\epsilon_N
 \log\frac{\det H_N[\tau_d]}{\det H_N[\sigma]}
 =4dq\log q+(8d\log2-4d)q+O_d(\log q).              \tag{OR38}
\]
For the low relation, rank zero contributes \(0\); rank one
has norm ratio between \([1+4(q+1)^2]^{-d}\) and \(1\), by
the same Jensen estimate applied to the full \(Q\).
It contributes \(O_d(\log q)\).
Using (OR36) at the two high relation ranks \(q,q+1\),
and the minus sign of the relation determinant in (OR2),
the net relation contribution to the baseline is
\[
 -4dq\log q-2dq\mathcal L(2,\pi)
 +O_h(\sqrt{q\log(q+2)}).                            \tag{OR39}
\]
Adding (OR38) and (OR39) proves
\[
 \BB[\tau_d]-\BB[\sigma]
 =-d[2\mathcal L(2,\pi)-8\log2+4]q
 +O_h(\sqrt{q\log(q+2)}).                            \tag{OR40}
\]
Together with (OR27), this proves (OR4)--(OR5).

\section{Integration scope}
All constants denoted \(C_h,c_h\) are fixed-source constants with
existence proved by the indicated estimates.  They have not been
evaluated numerically for an unspecified off-critical quartet.
The cofinal assertion retains the original packet cutoffs and uses,
in addition, the following literal finite inequalities wherever the
corresponding step is applied:
\[
 \begin{gathered}
 q\geq2^{16},\quad Y=2^{-16}q,\quad R\leq Y/4,\quad
 R^2/q^{7/4}\leq(5/4)2^{-28},\\
 K=q^{9/10}\geq kb,\quad
 T^{19/20}\geq2(k-1)b,\quad
 2T^{19/20}<T,\quad G\log T<\sqrt T/2,\\
 \epsilon_k\leq1/2,\quad
 \epsilon_{k-3}\leq e^{-1/2}/2,\quad k\geq V/b^2,\quad
 r\leq2q+1,\\
 \epsilon=\sqrt{\log(q+2)/q},\quad
 d/q<\epsilon/2,\quad 2\epsilon<1/8 ,
 \qquad q^{9/10}\leq T\leq128q .
 \end{gathered}                                      \tag{OR41}
\]
Here \(\epsilon_n\) is the original relative local-limit error
from the localization provider and \(r\) is the selected Taylor
rank in (OR22).  The Chernoff step also takes its explicit eventual
range \(C_hk/T+C_hk\log T/T^{19/20}\leq1/4\), enlarging the
fixed-source constants if needed.  The root-comparison appendix
retains its displayed exponentially small \(\varepsilon_q\leq1/2\).
Every guard holds eventually on each original fixed-multiplicity
sequence: \(q\asymp_h k^2\) for \(m=1\), \(q\asymp_h k^3\) for
\(m>1\), and \(R=O_h(k)\).  Thus these are finite-domain
specifications, not new assumptions replacing any missing estimate.

The fixed multiplicity, every root and repetition, both real
half-lines, the source-to-remainder map and its complete relation
Gram were used throughout.  The scalar \(A_k\) cancelled only at
the full quotient receiver (OR1).  The estimated complete
quotient determinant does not determine the seven absolute
allocations or projected arithmetic signs.  Those quantities
retain their original restricted metrics and same-class pairings.
The exact metric maps established here are (OR2), the Taylor
short sequence used in (OR23), and multiplication by \((y+i)^d\)
used in (OR27).

\begin{thebibliography}{9}
\bibitem{Ivic} A. Ivi\'c,
\emph{On the mean square of the Riemann zeta-function in short intervals},
arXiv:0803.0132v3, Section 2, (2.1) and the following \(E\)-analogue.
\url{https://arxiv.org/html/0803.0132}.
\bibitem{KVJ} H. T. Koelink and J. Van der Jeugt,
\emph{Bilinear generating functions for orthogonal polynomials},
arXiv:q-alg/9704016, Proposition 2.1, (2.9).
\url{https://arxiv.org/html/q-alg/9704016}.
\bibitem{DLMFhyper} NIST Digital Library of Mathematical Functions,
15.6.1, Euler integral for the regularized hypergeometric function.
\url{https://dlmf.nist.gov/15.6.E1}.
\bibitem{DLMFgamma} NIST Digital Library of Mathematical Functions,
5.11, Gamma asymptotic expansions.
\url{https://dlmf.nist.gov/5.11}.
\end{thebibliography}
\end{document}
